\documentclass[10pt, compress]{beamer}

\usetheme{Madrid}
\usecolortheme{default}
\setbeamertemplate{navigation symbols}{}
\setbeamertemplate{itemize items}[circle]

\usepackage[utf8]{inputenc}
\usepackage[T1]{fontenc}
\usepackage[french]{babel}
\usepackage{amsmath,amssymb}
\usepackage{graphicx}
\usepackage{booktabs}
\usepackage{subcaption}

\graphicspath{{figures/}}

% Macros
\newcommand{\Ltwo}{L^2}
\newcommand{\Dp}{D_p}
\newcommand{\Dw}{D_w}
\newcommand{\DK}{D_K}

\begin{document}

% ==============================================================================
\begin{frame}
  \title{Classification non-supervisée de données mixtes}
  \author{Pierre Chambet}
  \institute{Master TRIED \\ \vspace{0.5em} {\small Projet long proposé par V. Audigier, F. Bouhadjera et N. Niang \\ CNAM, Laboratoire CEDRIC, équipe MSDMA}}
  \date{Soutenance de Projet Long}
  \titlepage
\end{frame}

% ==============================================================================
\begin{frame}{1. L'Obstacle Mathématique et les Distances}
\textbf{Le Problème :}\\
Regrouper des individus caractérisés par une courbe $X(t) \in \Ltwo$ et un vecteur $Z \in \mathbb{R}^p$.

\vspace{1em}
Pour y parvenir, on commence par extraire formellement des dissimilarités séparées :
\begin{itemize}
    \item $D_0$ : Distance entre courbes (niveau)
    \item $D_1$ : Distance entre dérivées des courbes (forme)
    \item $D_s$ : Distance vectorielle classique sur l'espace $\mathbb{R}^p$
\end{itemize}
\end{frame}

% ==============================================================================
\begin{frame}{2. Les Stratégies de Surface}

Comment l'état de l'art tente de fusionner les distances mixtes ?

\vspace{1em}
\textbf{Stratégie A (Concaténation naïve) :} \\
On extrait les scores analytiques des courbes (FPCA $\rightarrow \eta$), et on les concatène brutalement avec le vecteur $Z$. 

\vspace{1em}
\textbf{Stratégie B (Somme pondérée $D_w$) :} \\
Combinaison euclidienne emboîtée des distances préalablement normalisées :
$$ \Dw(\alpha,\omega) = \sqrt{\omega \underbrace{\bigl[(1-\alpha)\tilde{D}_0^2 + \alpha\tilde{D}_1^2\bigr]}_{D_p(\alpha)^2} + (1-\omega)\tilde{D}_s^2} $$

\vspace{1em}
\textbf{Stratégie C (Produit de Noyaux $D_K$) :} \\
Multiplication des probabilités de similarité plutôt que somme de distances :\\
$K = K_f \cdot K_s$, où $K_f = \exp(-D_p^2 / 2\sigma_f^2)$.
$D_K$ est la distance :
$$D_K(i,j) = \sqrt{K(i,i) + K(j,j) - 2K(i,j)}$$
\end{frame}

% ==============================================================================
\begin{frame}{3. La Limite Structurelle Fondamentale}

\vspace{2em}
\begin{alertblock}{Le Verrou de l'Hybridation Tardive}
Que l'on additionne, concatène ou multiplie, on fusionne des informations de surface \textbf{après} les avoir calculées indépendamment. 
\end{alertblock}

\vspace{1em}
\textbf{Conséquence inévitable :}\\
On interdit mathématiquement à la courbe d'influencer analytiquement le vecteur. On omet de modéliser et de traiter la \textbf{covariance croisée} inter-modalités en amont.
\end{frame}

% ==============================================================================
\begin{frame}{4. La Solution : ACP Hybride (HFV)}
\textbf{Méthode :}\\
Diagonaliser une structure jointe en amont du clustering.

\begin{enumerate}
    \item \textbf{Réduction séparée} : Scores FPCA ($\eta$) et ACP classique ($\gamma$).
    \item \textbf{Ratio d'équilibrage $r$} : Rapport des variances.
    \item \textbf{Covariance Jointe $V$} :
    \[
      V = \begin{bmatrix} V_x & \mathbf{V_{xz}} \\ \mathbf{V_{zx}} & V_z \end{bmatrix}
    \]
    {\small $x$ = bloc fonctionnel (courbes), \quad $z$ = bloc vectoriel ($Z$)}
\end{enumerate}

\vspace{0.5em}
Les vecteurs propres de $V$ mélangent \textbf{par construction} composantes analytiques et coefficients discrets, autorisant une reconstruction informée dans le sous-espace $L^2$.
\end{frame}

% ==============================================================================
\begin{frame}{5. Le Matériel d'Évaluation (Données simulées et réelles)}
Nous utilisons des données labellisées.

\vspace{0.5em}

\textbf{1. Les Données Simulées (4 Scénarios)} \\
Simulation du niveau des courbes ($D_0$), de la forme ($D_1$) et du vecteur $Z$ :
\begin{itemize}
    \item \textbf{S1} : Signal fort et covariance mixte.
    \item \textbf{S2} : Bruit vectoriel fort.
    \item \textbf{S3} : Bruit fonctionnel fort.
    \item \textbf{S4} : Bruit dominant.
\end{itemize}
\vspace{0.5em}
\textbf{2. Les datasets réels} \\
\begin{table}
\centering
\renewcommand{\arraystretch}{1.1}
\resizebox{1\textwidth}{!}{
\begin{tabular}{l p{3.4cm} p{3.4cm} l}
\toprule
\textbf{Dataset} & \textbf{Courbe $X(t)$} & \textbf{Vecteur $Z$} & \textbf{Cible ($K$)} \\
\midrule
\textbf{Canadian Weather} & Température annuelle & Latitude, Long., Précip. & \textbf{4 régions} \\
\textbf{Berkeley Growth} & Courbe de croissance & Taille finale, poids & \textbf{2 sexes} \\
\textbf{Tecator} & Spectre d'absorbance & Eau, protéines & \textbf{3 taux de gras} \\
\bottomrule
\end{tabular}}
\end{table}
\end{frame}

% ==============================================================================
\begin{frame}{6. Le Problème de la Silhouette}

Les méthodes non-supervisées optimisent leurs paramètres ($\alpha$, $\omega$) en aveugle grâce à la \textbf{Silhouette moyenne}.

\vspace{1em}
\begin{alertblock}{Le Paradoxe Observé}
Le couple $(\alpha, \omega)$ maximisant la Silhouette s'écarte systématiquement de l'optimum maximisant la véritable structure physique (ARI terrain).
\end{alertblock}

\vspace{0.5em}
\textbf{Interprétation :}\\
Pour gonfler mathématiquement son score, la silhouette privilégie aveuglément des amas denses et sphériques fragmentant ainsi les structures de formes variables.
\end{frame}

% ==============================================================================
\begin{frame}{7. Preuve Empirique du Problème}

Évaluation contrefactuelle sur la Stratégie B : perte d'ARI causée par l'optimisation aveugle (Silhouette) vs optimisation Oracle (ARI).

\vspace{1em}
\begin{table}
\centering
\renewcommand{\arraystretch}{1.3}
\resizebox{0.95\textwidth}{!}{
\begin{tabular}{lccc}
\toprule
\textbf{Jeu de données} & \textbf{opt-Silhouette} & \textbf{opt-ARI} & \textbf{Chute ($\Delta$ ARI)} \\
\midrule
\textbf{Canadian Weather} & $0{.}616$ & $0{.}898$ & \textbf{${-} 0{.}283$} \\
\textbf{Berkeley Growth} & $0{.}424$ & $0{.}756$ & \textbf{${-} 0{.}332$} \\
\textbf{Tecator} & $0{.}151$ & $0{.}626$ & \textbf{${-} 0{.}475$} \\
\bottomrule
\end{tabular}}
\end{table}

\vspace{0.5em}
\textbf{Conséquence :} L'architecture non-supervisée s'effondre systématiquement en l'absence de compacité quasi-parfaite.
\end{frame}

% ==============================================================================
\begin{frame}{8. L'Instabilité Bootstrap (Canadian Weather)}

Face à la défaillance de la silhouette, la sécurisation du modèle passe obligatoirement par \textbf{l'Instabilité Bootstrap de Fang \& Wang} (Moins d'instabilité = rouge $\rightarrow$ bleu).

\vspace{1em}
\begin{figure}
  \centering
  \includegraphics[width=0.75\textwidth]{nselectboot_heatmap_canadian.png}
  \caption{Canadian Weather (4 Régions)}
\end{figure}

Les zones idéales $(\alpha, \omega)$ (croix bleues) fluctuent massivement selon le jeu de données.
\end{frame}

% ==============================================================================
\begin{frame}{9. Instabilité Bootstrap (Growth \& Tecator)}


\vspace{1em}
\begin{figure}
  \centering
  \begin{subfigure}{0.48\textwidth}
    \includegraphics[width=\linewidth]{nselectboot_heatmap_growth.png}\caption{Berkeley Growth (Sexe)}
  \end{subfigure}
  \hfill
  \begin{subfigure}{0.48\textwidth}
    \includegraphics[width=\linewidth]{nselectboot_heatmap_tecator.png}\caption{Tecator (Gras)}
  \end{subfigure}
\end{figure}
\end{frame}

% ==============================================================================
\begin{frame}{10. Instabilité Bootstrap (Simulations S1 \& S2)}

Evaluation de Fang \& Wang sur nos générateurs de simulation :
\vspace{1em}
\begin{figure}
  \centering
  \begin{subfigure}{0.48\textwidth}
    \includegraphics[width=\linewidth]{nselectboot_heatmap_sim_S1_seed1.png}\caption{S1}
  \end{subfigure}
  \hfill
  \begin{subfigure}{0.48\textwidth}
    \includegraphics[width=\linewidth]{nselectboot_heatmap_sim_S2_seed1.png}\caption{S2 (Signal fonctionnel fort)}
  \end{subfigure}
\end{figure}

L'invariant structurel est retrouvé par l'instabilité (la zone bleue de faible instabilité s'aligne bien avec le vrai nombre de classes $K=3$).
\end{frame}

% ==============================================================================
\begin{frame}{11. Instabilité Bootstrap (Simulations S3 \& S4)}

Dès qu'on approche du scénario S4 (bruit total dominant), l'invariant structurel s'effondre :
\vspace{1em}
\begin{figure}
  \centering
  \begin{subfigure}{0.48\textwidth}
    \includegraphics[width=\linewidth]{nselectboot_heatmap_sim_S3_seed1.png}\caption{S3 (Signal vectoriel fort)}
  \end{subfigure}
  \hfill
  \begin{subfigure}{0.48\textwidth}
    \includegraphics[width=\linewidth]{nselectboot_heatmap_sim_S4_seed1.png}\caption{S4 (Bruit total)}
  \end{subfigure}
\end{figure}

\end{frame}

% ==============================================================================
\begin{frame}{12. Résultats Globaux (Scores ARI Moyens)}

Clustering sur les données simulées :

\begin{table}
\centering
\resizebox{0.95\textwidth}{!}{
\begin{tabular}{lccccc}
\toprule
\textbf{Méthode} & \textbf{S1} & \textbf{S2} & \textbf{S3} & \textbf{S4} & \textbf{Moyenne} \\
\midrule
Stratégie A (FPCA$+Z$)        & $0{.}782$ & $0{.}524$ & $0{.}623$ & $0{.}315$ & $0{.}561$ \\
Stratégie B ($\Dw$)           & $0{.}873$ & $0{.}873$ & $0{.}239$ & $0{.}238$ & $0{.}556$ \\
Stratégie C ($\DK$)           & $\mathbf{0{.}938}$ & $\mathbf{0{.}887}$ & $0{.}464$ & $0{.}317$ & $\mathbf{0{.}651}$ \\
\midrule
HFV + $\DK$                   & $0{.}897$ & $0{.}802$ & $0{.}433$ & $0{.}301$ & $0{.}608$ \\
\midrule
Baseline $D_0$ (Niveau seul)  & $0{.}830$ & $0{.}830$ & $0{.}226$ & $0{.}226$ & $0{.}528$ \\
Baseline $D_s$ (Vecteur seul) & $0{.}567$ & $0{.}333$ & $\mathbf{0{.}511}$ & $0{.}247$ & $0{.}414$ \\
\bottomrule
\end{tabular}}
\end{table}

\textbf{Constats :}
\begin{itemize}
    \item HFV et Stratégie C mènent dès que le signal mixte est présent (S1, S2).
    \item L'hybridation échoue à restaurer les classes si le signal faiblit uniformément (S4).
\end{itemize}
\end{frame}

% ==============================================================================
\begin{frame}{13. Résultats sur les Jeux Réels Etiquetés}

Clustering sur les données réelles :
\begin{table}
\centering
\resizebox{1\textwidth}{!}{
\begin{tabular}{l ccc c cc}
\toprule
\textbf{Jeu} & \textbf{Strat.A} & \textbf{Strat.B} & \textbf{Strat.C} & \textbf{HFV} & \textbf{Base $D_0$} & \textbf{Base $D_s$} \\
\midrule
Canadian (Régions) & $\mathbf{0{.}748}$ & $0{.}616$ & $0{.}686$ & $0{.}686$ & $0{.}229$ & $0{.}616$ \\
Growth (Sexe)      & $\mathbf{0{.}682}$ & $0{.}424$ & $0{.}368$ & $0{.}368$ & $0{.}115$ & $0{.}424$ \\
Tecator (Gras)     & $0{.}439$ & $0{.}151$ & $0{.}371$ & $0{.}374$ & $0{.}151$ & $\mathbf{0{.}546}$ \\
\bottomrule
\end{tabular}}
\end{table}

\textbf{La Leçon de Tecator :}
Ici, seule la composition chimique discrète $D_s$ sépare les individus. Forcer l'hybridation avec la courbe d'absorbance \emph{nuit} au résultat.
\end{frame}

% ==============================================================================
\begin{frame}{14. Synthèse}

\begin{block}{Validation d'Architecture}
L'ACP Hybride (HFV) apporte une solution théorique robuste à la liaison $V_{yx}$ entre courbes et vecteurs, tenant tête empiriquement aux noyaux $D_K$ classiques.
\end{block}

\vspace{0.5em}
\begin{alertblock}{La limite fondamentale du pur Non-Supervisé}
La Silhouette dévaste trop souvent les géométries variables pour maximiser des "masses rondes" irréalistes.
\end{alertblock}

\end{frame}

% ==============================================================================
\end{document}
